% 第3章3.4.3退化换基（verify.py 验算）。三点算例但 u_st=1：T={s->a,a->t}，U={s->t}，x=(2,2,1)，费用8，
% pi=(2,1,0)，c̄_st=4-2+0=2>0 且在上界，反向进入；环上剩余：s->t 降1，s->a 增1，a->t 增0，步长0。
% a->t 出基进入 U，T={s->a,s->t}，流不变，pi=(4,3,0)，c̄_at=-2<=0，停止。
\begin{tikzpicture}[>=Stealth,line width=.8pt,font=\fontsize{8}{10}\selectfont,
  v/.style={circle,draw=FigureStroke,fill=FigureBlueFill,line width=.8pt,minimum size=6.5mm,inner sep=0pt,text=FigureInk},
  rv/.style={v,fill=FigureYellowFill},
  tr/.style={->,draw=FigureBlue,line width=1.4pt,shorten >=.5pt},
  lo/.style={->,draw=FigureMuted,line width=.8pt,dashed,shorten >=.5pt},
  up/.style={->,draw=FigureMuted,line width=.9pt,densely dotted,shorten >=.5pt},
  lab/.style={text=FigureInk,inner sep=1.5pt},
  pot/.style={text=FigureGreen},
  note/.style={text=FigureMuted,font=\FigureNoteFont\fontsize{7.5}{9}\selectfont}]
  \foreach \dx/\k in {0/A,5.2/B}{
  \begin{scope}[xshift=\dx cm]
  \node[v] (s\k) at (0,0) {$s$};
  \node[v] (a\k) at (1.7,1.4) {$a$};
  \node[v] (t\k) at (3.4,0) {$t$};
  \end{scope}}
  \draw[tr] (sA)--node[lab,above left=-1pt,align=right]{$2/3$\\$+$\,剩$1$} (aA);
  \draw[tr,draw=FigureRed] (aA)--node[lab,above right=-1pt,align=left]{$2/2$\\$+$\,剩$0$} (tA);
  \draw[up,draw=FigureGreen,line width=1.1pt] (sA)--node[lab,below]{$1/1$\,$U$\quad$-$\,剩$1$} (tA);
  \node[pot,below=1pt] at (sA.south) {$\pi=2$};
  \node[pot,above=1pt] at (aA.north) {$\pi=1$};
  \node[pot,below=1pt] at (tA.south) {$\pi=0$};
  \node[note] at (1.7,-.95) {$\bar c_{s,t}=2>0$，$\theta=0$};
  \node at (1.7,-1.5) {(a) 树弧$a\to t$已在上界};
  \draw[tr] (sB)--node[lab,above left=-1pt]{$2/3$} (aB);
  \draw[up] (aB)--node[lab,above right=-1pt]{$2/2$\,$U$} (tB);
  \draw[tr] (sB)--node[lab,below]{$1/1$} (tB);
  \node[pot,below=1pt] at (sB.south) {$\pi=4$};
  \node[pot,above=1pt] at (aB.north) {$\pi=3$};
  \node[pot,below=1pt] at (tB.south) {$\pi=0$};
  \node[note] at (6.9,-.95) {$\bar c_{a,t}=-2\leq0$};
  \node at (6.9,-1.5) {(b) 流不变，树基改变};
\end{tikzpicture}
